\documentclass[10pt,a4paper]{amsart}
\usepackage[margin=19mm]{geometry}
\usepackage{amsmath,amssymb,mathtools}
\usepackage[scr=rsfs,cal=euler]{mathalfa}
\usepackage{xfrac}
\usepackage[unicode,hidelinks,bookmarks=false]{hyperref}
\newcommand{\ie}{i.e.\ }
\newcommand{\cf}{cf.\ }
\newcommand{\resp}{respectively\ }
\newcommand{\Subsection}{\subsection}
\providecommand{\nobreakdash}{\nobreak\hbox{-}}
\setlength{\emergencystretch}{2em}
\multlinegap0pt
\usepackage{bm}
\usepackage{bbm}
\usepackage{dsfont}
\usepackage{xspace}
\usepackage{cancel}
\usepackage{mathtools}
\usepackage[shortlabels]{enumitem}
\usepackage{amsthm}
\makeatletter
\@ifundefined{theorem}{\newtheorem{theorem}{Theorem}}{}
\@ifundefined{lemma}{\newtheorem{lemma}[theorem]{Lemma}}{}
\makeatother
\newcommand{\PP}{\mathcal{P}}
\newcommand{\R}{\mathbb R}
\newcommand{\conv}{\mathrm{conv}\,}
\newcommand{\iconv}{\mathrm{iconv}\,}
\newcommand{\lc}{\preceq_{\textnormal{c}}}
\newcommand{\lic}{\preceq_{icx}}
\title[Weak Optimal Transport: When is the Dual Potential Convex?]{Weak Optimal Transport: When is the Dual Potential Convex?}
\author{Filip Pramenkovi\'c}
\date{Source: arXiv:2507.07200v2 (CC BY 4.0)}
\begin{document}
\maketitle

\noindent\emph{Source and licence.} Weak Optimal Transport: When is the Dual Potential Convex?. Version arXiv:2507.07200v2 (CC BY 4.0), distributed under the Creative Commons Attribution 4.0 International licence, as recorded in the arXiv licence metadata for this exact version and shown on the abstract page https://arxiv.org/abs/2507.07200v2. Full rights notice: licenses/arxiv-2507.07200-CC-BY-4.0.txt.\par\smallskip

\noindent\textit{Editorial scope.} Counted unit: the Lemma whose source label is \texttt{rc stabil lemma} in the article (source0.tex lines 421--424, source label \texttt{rc stabil lemma}), delivered together with the article's own complete original proof of it (source0.tex lines 426--465, one \texttt{proof} environment of 40 lines). No mathematics is written, repaired or re-derived: statement and proof are the article's own text line for line. Required context retained uncounted: repr lemm (lines 321--334); conv repr bullet (lines 323--332); prop1 (lines 323--332); prop4 (lines 323--332); prop6 (lines 323--332). Every definition, display and label of those lines is retained unchanged. Omitted from this package and not redistributed: 1--320, 333--420, 425--425, 466--1007. Nothing inside a retained excerpt is omitted or altered: every retained body is delivered exactly as the article's lines are written, so no omission receipt of any kind accompanies this package. Citations: the retained excerpts carry no citation command at all, so no bibliography entry of the article is transcribed and no reference is redistributed. Reference adaptation: none was needed and none was made; every reference inside the delivered unit resolves inside it, so no unresolved reference or citation is produced. Printed slips kept literal: no printed word, symbol, hypothesis or number of the counted statement or of its proof was altered. Layout and class: the article's own class is \texttt{amsart} with packages including graphicx, hyperref, inputenc, babel, fontenc, datetime, amsfonts, bbm, amssymb, amsthm, xcolor, mathtools, accents, enumitem; the wrapper is the lane's pinned \texttt{amsart} editorial wrapper at 10pt on A4, which restates the theorem-like environments the retained text needs and adds the article's own macro definitions that the retained text needs (\texttt{\textbackslash PP}, \texttt{\textbackslash R}, \texttt{\textbackslash conv}, \texttt{\textbackslash iconv}, \texttt{\textbackslash lc}, \texttt{\textbackslash lic}), with extra wrapper packages (bm, bbm, dsfont, xspace, cancel, mathtools, enumitem). The delivered numbering is the unit's own. Assets: no retained span contains a figure environment, a graphics-inclusion command, a PDF-page inclusion, a table or a TikZ picture; no image file is copied into this package, the package declares no asset dependency and no image file is redistributed with it. Licence: version arXiv:2507.07200v2 of this article is distributed under the Creative Commons Attribution 4.0 International licence, as recorded in the arXiv licence metadata for this exact version and shown on the abstract page https://arxiv.org/abs/2507.07200v2. A full rights notice is delivered at licenses/arxiv-2507.07200-CC-BY-4.0.txt. Characters: every retained excerpt is pure ASCII, so the delivered engine, XeLaTeX with its default Latin Modern text font, reports no missing character.
\begin{lemma} \label{repr lemm}
    Let $\psi \in {\mathcal{C}}_{b,p}(\R^d)$. Then we have 
    \begin{enumerate}[label=(\roman*)]\label{repr enac}
     \item \label{conv repr bullet}  $\conv{\psi}(y)=\inf\left\{\rho(\psi) \mid \delta_y\lc\rho,\; \rho\in\PP_p(\R^d)\right\},$
     \item \label{iconv repr bullet}  $\iconv{\psi}(y)=\inf\left\{\rho(\psi) \mid \delta_y\lic\rho,\; \rho\in\PP_p(\R^d)\right\},$
      \item \label{prop1} $\conv{\psi}\leq \psi$,
     \item \label{prop3}$\psi_1,\psi_2\in{\mathcal{C}}_{b,p}(\R^d)$ and $\psi_1\leq\psi_2$, then $\conv{\psi_1}\leq \conv{\psi_2},$
      % \item \label{prop5}if $\psi $ is convex, then $\conv{\psi}=\psi$,
  %  \item \label{prop2}$\widehat{\psi}$ is real-valued,
    \item \label{prop4}$\conv{\psi}\in{\mathcal{C}}_{b,p}(\R^d),$
    \item \label{prop6}$\conv{\psi_n}\downarrow \conv{\psi}$, as $n\to\infty$, where $\psi_n:=\psi+\frac{1}{n}|\cdot|^p.$
    \end{enumerate}
Additionally, properties \ref{prop1}--\ref{prop6} also hold for the increasing convex hull.
\end{lemma}

\begin{lemma} \label{rc stabil lemma}
    Let $C:X\times\PP_p(\R^d)\to[0,\infty]$ be measurable and $\psi\in {\mathcal{C}}_{b,p}(\R^d)$. If $C$ is $\lc$-decreasing in the second argument, then we have $$ \psi^C=({\conv{\psi}})^C,$$
    and if $C$ is $\lic$-decreasing in the second argument, then we have $$ \psi^C=({\iconv{\psi}})^C.$$
\end{lemma}

\begin{proof} We prove the first claim, as the second is proven identically.
    Fix $\psi\in{\mathcal{C}}_{b,p}(\R^d)$. Then
    $$ \psi^C(x)=\inf_{\rho\in\PP_p(\R^d)} \rho(\psi)+C(x,\rho)\geq \inf_{\rho\in\PP_p(\R^d)} \rho(\conv\psi)+C(x,\rho)=(\conv{\psi})^C(x),$$
    where the inequality is by Lemma \ref{repr lemm}\textit{\ref{prop1}}. 
    
    The converse inequality is equivalent to the statement
   \begin{equation} \label{nejjed}
        \forall\varepsilon>0\; \forall\rho\in\PP_p(\R^d) \;\exists\widetilde{\rho}\in\PP_p(\R^d):\widetilde{\rho}(\psi)+C(x,\widetilde{\rho})\leq {\rho}(\conv{\psi})+C(x,\rho)+\varepsilon.
   \end{equation}
    Fix $\varepsilon>0$ and $\rho\in\PP_p(\R^d).$ Then by Lemma \ref{repr lemm}\textit{\ref{prop6}} we have
    $\conv{\psi_n}\downarrow \conv{\psi}$, as $n\to\infty$, where $\psi_n:=\psi+\frac{1}{n}|\cdot|^p$. Since $\rho(\conv{\psi_1})<\infty$, we have by monotone convergence 
    $$\rho(\conv{\psi_n})\xrightarrow{n\to\infty}\rho(\conv{\psi}).$$
    So fix $N\in\mathbb N$ sufficiently large such that
   \begin{equation} \label{suff small}
        \rho(\conv{\psi_N})-\rho(\conv \psi)\leq \varepsilon/2.
   \end{equation}
    Now by Lemma \ref{repr lemm}\textit{\ref{conv repr bullet}} applied to $\psi_N$ we have that the set 
    $$A:=\left\{(y,\xi)\in \R^d\times\PP_p(\R^d) \mid \conv{\psi_N}(y)+\varepsilon/2\geq \xi(\psi_N), \;\delta_y \lc\xi\right\}$$
    satisfies $\text{proj}_1A=\R^d$. It is straightforward to verify that $A$ is Borel measurable and by the von Neumann uniformization theorem, we may extract a universally measurable selection $\{\xi_y\}_{y\in \R^d}\subseteq\PP_p(\R^d)$, such that
    \begin{equation} \label{selec prop}
        \conv{\psi_N}(y)+\varepsilon/2\geq \xi_y(\psi_N), \;\delta_y \lc\xi_y.
    \end{equation}
   We define $\widetilde{\rho}$ by $d\widetilde{\rho}(z):=\int_{\R^d} d\xi_y(z)d\rho(y).$ Integrating the inequality in (\ref{selec prop}) with respect to $\rho$ yields
   \begin{equation} \label{bitna}
       \rho(\conv{\psi_N})+\varepsilon/2\geq \widetilde{\rho}(\psi_N).
   \end{equation}
   On the one hand, since $\psi$ is bounded from below we have
$$\widetilde{\rho}(\psi_N)\geq \ell +1/N\widetilde{\rho}(|\cdot|^p),$$
and on the other hand, by Lemma \ref{repr lemm}\textit{\ref{prop4}}, 
$\rho(\conv{\psi_N}) <\infty.$
These two inequalities combined with (\ref{bitna}) yield that $\widetilde{\rho}\in\PP_p(\R^d)$.

Now notice that (\ref{suff small}) and (\ref{bitna}) imply
\begin{equation} \label{bistna}
       \rho(\conv{\psi})+\varepsilon\geq \widetilde{\rho}(\psi_N)\geq\widetilde{\rho}(\psi).
   \end{equation}
This partially shows the inequality in (\ref{nejjed}). To finalize the proof, note that for arbitrary convex $f$, we have
$$\widetilde{\rho}(f)=\int \xi_y(f)d\rho(y)\geq\int \delta_y(f)d\rho(y)=\rho(f),$$
i.e., $\rho\lc\widetilde{\rho}$, which by the monotonicity assumption implies $C(x,\rho)\geq C(x,\widetilde{\rho})$.
\end{proof}

\end{document}
